\begin{lemma}[The corresponding terminal series tree is the selected closure]
Let \(T\) be an extended Tashkinov tree of positive rung number \(n\),
presented by a certified series \(\mathcal T\), and let \(T_n\) denote its
terminal ladder tree.  Fix a coloring \(\sigma_n\), tree-sequences \(S\) and
\(T^*\), and a \noderef{CorrespondingExtendedTashkinovTree} certificate
\(\mathcal C\) exhibiting \(T^*\) as corresponding to \(T\) under
\(\sigma_n\), with replacement series \(\mathcal T^*\).  Suppose that one
of the following exact seed conditions holds:
\begin{itemize}
\item the terminal kind in \(\mathcal T\) is PE and \(S=T_n\); or
\item the terminal kind is SE or RE, and there is an edge identity \(f\)
such that the terminal connecting-edge set is exactly \(\{f\}\) and
\noderef{TreeSequenceAdjoinEdge}\((T_n,f,S)\) holds.
\end{itemize}
If \(T^*\) is a \noderef{TashkinovClosure} of \(S\) under \(\sigma_n\),
then the final output tree \(T^*_{n+1}\) of \(\mathcal T^*\) is exactly
\(T^*\) as an ordered edge-valued tree-sequence.
\end{lemma}
\begin{proof}
Write \(\mathcal E^*\) for the replacement ETT carried by \(\mathcal C\).
By the rung-agreement and dependent-series-presentation fields in
\noderef{CorrespondingExtendedTashkinovTree}, we may identify the rung number
of \(\mathcal E^*\) with \(n\) and its certified series with
\(\mathcal T^*\).  Since \(n\geq 1\) by
\noderef{ExtendedTashkinovTree}, the terminal stored index is the successor
of the index \(n-1\), and its output-tree index is the following successor,
namely the final tree index \(n+1\).  The ladder-tree, connecting-edge, and
extension-kind agreement fields of \(\mathcal C\), evaluated at the terminal
stored index, give
\[
 T_n^*=T_n,\qquad F_n^*=F_n,\qquad \Theta_n^*=\Theta_n.       \tag{1}
\]
Also, the final-generating-coloring field gives
\(\varphi_n^*=\sigma_n\).                                      \tag{2}

By \noderef{TashkinovSeries}, the terminal stored index of
\(\mathcal T^*\) has a step certificate.  Its four possible forms are
listed in \noderef{TashkinovSeriesStep}.  The
initial form has no extension kind, whereas the hypotheses and (1) give
PE, SE, or RE, so that form is impossible.  Likewise, in each of the three
noninitial forms, (1) rules out the two forms whose recorded kind differs
from \(\Theta_n\).  It remains to handle the matching PE, SE, and RE forms.

Suppose first that \(\Theta_n=\mathrm{PE}\).  The PE terminal-step
certificate says that \(T^*_{n+1}\) is a closure, under
\(\varphi_n^*\), of its terminal ladder tree \(T_n^*\).  By (1), (2), and
the PE seed hypothesis, this says precisely that \(T^*_{n+1}\) is a
closure of \(S\) under \(\sigma_n\):
\[
              \operatorname{TashkinovClosure}
                (\sigma_n,S,T^*_{n+1}).                          \tag{3}
\]

Now suppose that \(\Theta_n\) is SE or RE.  Let \(f\) be the edge in the
seed hypothesis.  In either matching terminal-step certificate, there is
an exact edge identity \(f^*\), the stored connecting-edge set is the
singleton \(\{f^*\}\), an intermediate tree \(A^*\) satisfies
\noderef{TreeSequenceAdjoinEdge}\((T_n^*,f^*,A^*)\), and
\(T^*_{n+1}\) is a closure of \(A^*\) under \(\varphi_n^*\).  Equality
\(F_n^*=F_n\) in (1), together with the two singleton descriptions, yields
\(f^*=f\) as exact edge identities.  Equality \(T_n^*=T_n\) then shows
that both \(A^*\) and \(S\) are obtained by adjoining the same exact edge
\(f\) to the same ordered tree-sequence \(T_n\).  Therefore
\noderef{TreeSequenceAdjoinEdgeUnique} gives \(A^*=S\).  Substituting this
equality and (2) into the terminal-step closure clause again proves (3).

The assumed closure of \(S\) is \(T^*\), while (3) shows that another
closure of the same seed under the same coloring is \(T^*_{n+1}\).
Hence \noderef{TashkinovClosureVertexSetUnique} gives
\[
                         V(T^*)=V(T^*_{n+1}).                    \tag{4}
\]
Finally, the below-final field of the replacement
\noderef{ExtendedTashkinovTree} presentation says that \(T^*\) is an exact
ordered initial segment of \(T^*_{n+1}\).  Apply
\noderef{TreeSequenceInitialSegmentVertexSetEq} to that initial-segment
relation and (4).  It identifies the full ordered vertex and exact-edge
data, and therefore yields \(T^*=T^*_{n+1}\), as required.
\end{proof}
